\chapter*{Lời mở đầu}
\addcontentsline{toc}{chapter}{Lời mở đầu}

Một giao diện có bong bóng tin nhắn chưa chắc đã là một chatbot có trí nhớ. Model ngôn ngữ chỉ
thấy dữ liệu nằm trong request hiện tại; nếu ứng dụng không gửi lại thông tin cũ, câu hỏi
``Tôi tên gì?'' không có cách nào liên hệ với câu ``Tôi tên Nhật'' ở request trước. Từ quan sát
nhỏ này xuất hiện cả một lớp bài toán hệ thống: lưu message ở đâu, phân tách nhiều cuộc trò
chuyện thế nào, đưa bao nhiêu lịch sử vào context, tóm tắt phần cũ ra sao, và thông tin nào được
phép sống lâu hơn một thread.

Giáo trình này xây mô hình tinh thần trước rồi mới thêm framework. Học viên bắt đầu bằng một
danh sách Python thuần, chuyển nó thành state của LangGraph, thêm checkpointer, quan sát hai
\texttt{thread\_id} độc lập, rồi mới học trimming, running summary và long-term Store. Phần cuối
ghép các mảnh ấy vào FastAPI và thiết kế API cho một LLM Gateway nhiều provider.

\begin{aidefinition}[Phạm vi]
Cuốn sách tập trung vào bộ nhớ hội thoại và context engineering. Qdrant, RAG và agent gọi tool
chỉ được nhắc ở ranh giới cần phân biệt; chúng không phải điều kiện để hoàn thành các lab.
\end{aidefinition}

\begin{aiintuition}[Nguyên tắc xuyên suốt]
Model không sở hữu trí nhớ. Ứng dụng sở hữu dữ liệu; LangGraph giúp tổ chức state, checkpoint
và luồng cập nhật; model chỉ nhận một lát context do ứng dụng lựa chọn ở thời điểm gọi.
\end{aiintuition}

Mọi ví dụ cốt lõi dùng một ``model quyết định'' bằng quy tắc. Nhờ vậy học viên kiểm tra được
memory mà không phụ thuộc mạng, API key hay độ ngẫu nhiên của LLM. Khi cơ chế đã đúng, ta thay
model giả bằng DeepSeek hoặc Gemini mà không đổi invariant của hội thoại.

\section*{Đối tượng và điều kiện đầu vào}

Tài liệu dành cho người đã biết Python cơ bản, dictionary, list, function và cách chạy một ứng
dụng FastAPI; không giả định đã dùng LangChain hay LangGraph. Python 3.11 trở lên và công cụ
\texttt{uv} là cấu hình đề xuất. LangGraph hiện yêu cầu Python 3.10 trở lên và được cài độc lập
với các gói provider \parencite{langgraphInstall}.

\section*{Kết quả đầu ra quan sát được}

Sau khi hoàn thành, học viên có thể:
\begin{enumerate}
  \item giải thích vì sao UI history khác model context và persistent history;
  \item dựng graph \texttt{START -> call\_model -> END};
  \item dùng checkpointer và \texttt{thread\_id} để tách cuộc trò chuyện;
  \item chọn trim hoặc summary dựa trên chi phí và yêu cầu giữ thông tin;
  \item dùng \texttt{user\_id} và Store cho memory xuyên thread;
  \item thiết kế API chống trộn dữ liệu giữa người dùng;
  \item viết test chứng minh các invariant trên.
\end{enumerate}
